\begin{table}[t]
\tbl{\label{tab:diophant}}
{\begin{tcolorbox}[tabularx={p{3cm}|p{2cm}|p{6.3cm}|},enhanced,fonttitle=\bfseries\large,fontupper=\normalsize\sffamily,
colback=yellow!10!white,colframe=red!50!black,colbacktitle=blue!30!white,
coltitle=black,title=Creating Diophantine Systems]
\hline
Command & Arguments & Purpose \\
\hline
\hline
\verb'-label'  &  label  &  Use the given name as file name. \\
\hline
\verb'-coefficient_matrix'  &  A   &  Set the coefficient matrix to the previously created vector with format information. \\
\hline
\verb'-coefficient_matrix_csv'  &  fname   &  Read the coefficient matrix from the given csv-file.\\
\hline
\verb'-RHS'  &  options   &  Set RHS info for a set of $m$ rows. Possible types of conditions are EQ, LE, INT, ZOR. This command can be repeated. Example: mult=2,EQ=1 is condition for 2 rows where the RHS is equal to zero. Example: mult=1,LE=3 is one condition where the RHS is at most three. Example: mult=1,INT=1to2 is one condition where the RHS is between 1 and 2. Example: mult=1,ZOR=2 is one condition where the RHS is either 0 or 2. \\
\hline
\verb'-RHS_csv'  &  fname   &  Read the RHS information from the given csv file. \\
\hline
\verb'-RHS_constant'  &  options   &  Set the RHS as in \verb'-RHS' (without the mult component).\\
\hline
\verb'-x_max_global'  &  $a$   &  Set the upper bound for all variables to $a$ \\
\hline
\verb'-x_min_global'  &  $a$   &  Set the lower bound for all variables to $a$ \\
\hline
\verb'-x_bounds'  &  list-of-values   &  Set the lower and upper bounds for all variables. \\
\hline
\verb'-x_bounds_csv'  &  fname   &  Read the lower and upper bounds for all variables from the given file. \\
\hline
\verb'-has_sum'  &  $s$   &  Force the sum of the variables to be $s.$ \\
\hline
\verb'-problem_of_Steiner_type'  &  N fname  &  Create an instance of an exact cover problem based on the matrix in the given file. The underlying set has size $N.$ \\
\hline
\verb'-maximal_arc'  &  $s$ $d$ secants subset  &  Create system for a maximal arc of size $s$ and degree $d$ in $\PG(2,q)$. Use the given set of two pencil lines. The subset picks the lines from the given pencils which are external.  \\
\hline
\verb'-field'  &  F   &  Use the field F. Related to \verb'-maximal_arc'. \\
\hline
\end{tcolorbox}}
\end{table}
